\documentclass[../../main.tex]{subfiles}
\begin{document}

\begin{flushright}
\footnotesize\textit{【自動文字起こし・要確認】}
\end{flushright}

{\bf [解]}
$a^3+b^3=(a+b)(a^2-ab+b^2)$である．以下合同式の法を$3$とする．対称性から，$(a,b)\equiv(1,1),(1,2),(2,2)$の時をしらべれば良い．

\begin{itemize}
\item[$\circ$] $(a,b)\equiv(1,1)$の時，$a^3+b^3\equiv2$となり不適．
\item[$\circ$] $(a,b)\equiv(2,2)$の時，$a^3+b^3\equiv16\equiv1$，同様に不適．
\end{itemize}

だから，$(a,b)\equiv(1,2)$が必要である．この時，$a+b\equiv0$，$a^2-ab+b^2\equiv0$から，$a^3+b^3$は少なくとも$9$で割り切れる．そこで，$m,n\in\mathbb{Z}_{\ge0}$として，
\begin{align}
a=3m+1，\quad b=3n+2\quad\cdots①
\end{align}
とおく．
\begin{align}
a+b&=3(m+n+1)，\\
a^2-ab+b^2&=(9m^2+6m+1)+(9n^2+12n+4)-(9mn+3n+6m+2)\\
&=9m^2+9n^2-9mn+9n+3
\end{align}

から，$a^2-ab+b^2$は$9$で割り切れず，$a+b$が$81$で割り切れるには$a+b$が$27$で割り切れなければならない．そこで，$\ell\in\mathbb{N}$として
\begin{align}
a+b=27\ell\quad\cdots②
\end{align}
とする．それが条件．$A=a^2+b^2$として，
\begin{align}
A&=\dfrac{1}{2}\left\{(a+b)^2+(a-b)^2\right\}\\
&=\dfrac{1}{2}\left\{(27\ell)^2+(27\ell-2b)^2\right\}
\end{align}

$\ell=1$の時，これは$b=14(\equiv2)$で最小値$365$をとる．この時$a=13$である．一方，$\ell\ge2$の時
\begin{align}
A\ge\dfrac{1}{2}(27\ell)^2\ge\dfrac{1}{2}\cdot4\cdot27^2=729\cdot2>365
\end{align}
から，$A$が$365$より小さくなることはない．以上から，$(a,b)=(13,14)$で$\min(a^2+b^2)=365$．この時対称性から$(a,b)=(14,13)$も解である．

\bigskip
\noindent{\bf [解2]}
$(a+b)^3=a^3+b^3+3ab(a+b)$から，$a^3+b^3\equiv0\pmod{81}$の時，$a+b$が$3$で割り切れ
\begin{align}
a+b=3c
\end{align}
とおける．代入・整理して
\begin{align}
9abc=27c^3-(a^3+b^3)
\end{align}
だから，$a,b\not\equiv0\pmod3$から，$c\equiv0\pmod3$である．この時
\begin{align}
a^3+b^3=27c^3-9abc
\end{align}
より，$a^3+b^3$は$81$で割り切れる．よって，条件は，$\ell\in\mathbb{N}$として
\begin{align}
\ulcorner\, a,b\not\equiv0\pmod3\text{かつ}a+b=27\ell\,\urcorner
\end{align}
（以下省略）

\bigskip
\noindent{\bf [注]}
最後の$a+b=27\ell$以下，図示してみると$\ell=1$で$\min$なのは明らかで，$\ell\ge2$の時
\begin{align}
a^2+b^2\ge(27)^2+(27)^2>365
\end{align}
を得る．

\newpage

\end{document}
